
\medskip

%Un jeu est constitué de quatre familles de cartes:
%banane ; 
%prune ; 
%citron ;
%fraise.
%
%Voici la répartition des cartes de la famille banane.
%
%\begin{center}
%\begin{tabularx}{\linewidth}{|m{3.5cm}|*{5}{>{\centering \arraybackslash}X|}}\hline
%Nombre de banane(s)	&1&2&3&4&5\\ \hline
%Nombre de cartes	&5&3&3&2&1\\ \hline
%\end{tabularx}
%\end{center}
%
%La répartition est la même pour les cartes avec les autres fruits.

\medskip

\begin{enumerate}
\item %Montrer que ce jeu a $56$ cartes.
Il y a dans chaque famille : $5 + 3 + 3 + 2 + 1 = 14$ cartes, donc en tout $4 \times 14 = 56$~cartes.

%Joanne mélange toutes les cartes. Son frère Jack prend une carte au hasard. On admet que chaque carte a la même chance d'être choisie.

%Soit $P$ l'évènement: \og Jack obtient une carte de la famille prune \fg.
\item %Quelle est la probabilité de l'évènement $P$ ?
Il y a 14 cartes prune sur 56 cartes, donc $p(P) = \dfrac{14}{56} = \dfrac14 = 0,25$.
\item 
	\begin{enumerate}
		\item %Quel est l'évènement contraire de $P$ ?
L'évènement contraire de $P$ est \og Jack ne tire pas une carte de la famille prune \fg.
		\item %Quelle est la probabilité de l'évènement contraire de $P$ ?
		La probabilité est égale à $1 - \dfrac14 = \dfrac34 = 0,75$.
	\end{enumerate}	
\item %Quelle est la probabilité d'obtenir une carte avec quatre fruits ?
Il y a dans chaque famille 2 cartes avec 4 fruits, donc en tout 8 cartes avec 4 fruits.

La probabilité de tirer l'une de ces cartes est donc : $\dfrac{8}{56} = \dfrac{8 \times 1}{8 \times 7} = \dfrac{1}{7}$.
\end{enumerate}

\bigskip

